\documentclass[12pt]{article}
\usepackage{helvet}
\usepackage{graphicx}
\usepackage{geometry}
\geometry{
    a4paper,
    top = 2cm,
    left = 2cm,
    right = 2cm,
    bottom = 3cm
}
\usepackage[utf8]{inputenc}
\usepackage[spanish,mexico]{babel} % Renombra en automático los elementos a español
\addto\captionsspanish{
  \renewcommand{\figurename}{Imagen}
  \renewcommand{\tablename}{Tabla}
}
\usepackage{float}  % Etiqueta 'H' (mayus) en figuras (imágenes)
\usepackage{array}  % Padding en tablas
\usepackage{tabularx}
\usepackage{url}    % Etiqueta url en .bib
\usepackage{booktabs}   % Manejo de caudros (tablas)
\usepackage[T1]{fontenc}    % Codificación de fuentes
\usepackage[autostyle, spanish=mexican]{csquotes}       % Paquete para comillas tipográficas
\usepackage{comment}    % Comentar bloques de líneas mediante begin-end
\usepackage{multicol}   % Para poner contenido a doble columna
\usepackage{enumitem}   % Para utilizar [noitemsep, topsep=0pt] en itemize
\usepackage{amsmath}
\usepackage{mathabx} % Para testar en math mode
\usepackage{comment}
\usepackage{longtable}
\usepackage[hidelinks]{hyperref}
\usepackage{caption}
\usepackage[
  backend=biber,
  style=apa,
  sorting=nyt
]{biblatex}
\DeclareLanguageMapping{spanish}{spanish-apa}
\renewcommand{\familydefault}{\sfdefault}
\newenvironment{localsubgroup}
  {$\left\{\tabular{l}}
  {\endtabular\right.$}


\addbibresource{bibliography/Creacion_Curriculum_Vitae.bib}

\begin{comment}
INDICACIONES DEL PROFESOR:
Fecha de entrega: 25 de octubre de 2026
Título: Creación Currículum Vitae
Número: Sin número
Tipo: Tarea

Instrucciones del profesor:

Objetivo

Diseñar un currículum vitae que refleje de forma organizada la formación académica, habilidades e intereses, fortaleciendo la preparación para futuras oportunidades laborales.

Complementa con un documento word , las paginas donde se pueden enviar solicitudes de trabajo.  Algunos ejemplos mínimos 5 de envío de curriculum para solicitar trabajo.

DIRECTIVAS ESPECÍFICAS DEL ALUMNO:
- Documento complementario enfocado en las plataformas y páginas para envío de solicitudes de empleo e internships en ingeniería de software.
- Fuente central de vacantes: repositorio SimplifyJobs Summer 2027 Internships (https://github.com/SimplifyJobs/Summer2027-Internships/blob/dev/README.md#-software-engineering-internship-roles).
- Sin prefijos numéricos en nombres de archivo.
- Sin recursos de media.
- Extensión flexible requerida para cubrir los objetivos, sin sujeción al mínimo de 18 páginas.
- Inclusión de ejemplos estructurados de envío a partir del perfil real del alumno aplicado a ByteDance / TikTok.
- Conclusiones personales orientadas a la preparación temprana desde la universidad y la aplicación del bagaje estudiantil en procesos de reclutamiento técnico.

ÍNDICE APROBADO:
1. Directorio de plataformas de reclutamiento [~3.0 págs.]
   1.1 Portales corporativos propietarios [~1.5 págs.]
   1.2 Sistemas de seguimiento de candidatos [~1.5 págs.]
2. Casos prácticos de envío de solicitudes [~3.5 págs.]
   2.1 Solicitud en ByteDance y TikTok
   2.2 Solicitud en Amazon
   2.3 Solicitud en KnowBe4
   2.4 Solicitud en VSP Vision
   2.5 Solicitud en PTC
3. Conclusiones personales [~1.0 pág.]
Referencias [~0.5 pág.]
\end{comment}

\begin{document}
\def\UrlBreaks{\do\/\do-\do_\do.\do?\do&\do=\do0\do1\do2\do3\do4\do5\do6\do7\do8\do9\do\a\do\b\do\c\do\d\do\e\do\f\do\g\do\h\do\i\do\j\do\k\do\l\do\m\do\n\do\o\do\p\do\q\do\r\do\s\do\t\do\u\do\v\do\w\do\x\do\y\do\z\do\A\do\B\do\C\do\D\do\E\do\F\do\G\do\H\do\I\do\J\do\K\do\L\do\M\do\N\do\O\do\P\do\Q\do\R\do\S\do\T\do\U\do\V\do\W\do\X\do\Y\do\Z}
\pagenumbering{roman}

\begin{titlepage}
  \thispagestyle{empty}
  \newgeometry{left = 1cm, right = 1cm, top = 1cm, bottom = 2cm}

  \noindent
  \begin{minipage}[c]{3.5cm}
    \includegraphics[width=3.5cm]{media/titlepage/Escudo-UNAM.pdf}
  \end{minipage}\hfill
  \begin{minipage}[c][3.5cm][c]{\dimexpr\textwidth-7cm\relax}
    \centering
    {\large\textbf{Universidad Nacional Autónoma de México}\\[0.6em] 
    \Large\textbf{Facultad de Ingeniería}\\[0.6em]
    \large Sistemas Operativos}
  \end{minipage}\hfill
  \begin{minipage}[c]{3.5cm}
    \raggedleft
    \includegraphics[width=3.5cm]{media/titlepage/FI.png}
  \end{minipage}
\centering\huge
  \vfill
    Grupo 1 \\
    Ing. Angel Isidro Mercado
  \vfill
    \textbf{Tarea Creación Currículum Vitae}
  \vfill\large
    320304435 - Cabrera Rojas Oscar
  \vfill

  \noindent\footnotesize \hfill Ciudad Universitaria, CDMX a 25 de octubre del 2026
\end{titlepage}
\restoregeometry

\tableofcontents

\pagenumbering{arabic}
\setcounter{page}{1}

\section{Directorio de plataformas de reclutamiento}
\label{sec:directorioPlataformas}

El proceso para conseguir una pasantía de verano (\textit{internship}) o un primer empleo en la industria del software se apoya actualmente en una serie de herramientas digitales diseñadas para recibir, clasificar y evaluar miles de solicitudes. Las empresas del sector rara vez procesan las postulaciones revisando archivos de manera manual una por una; en su lugar, utilizan sistemas de seguimiento de candidatos (\textit{applicant tracking system}, ATS) y portales de empleo corporativos que extraen el texto de los currículums (\textit{parsing}), identifican datos clave y canalizan a los aspirantes hacia pruebas técnicas automatizadas.

Entender cómo funcionan estas plataformas y qué tipo de información solicitan resulta clave para estructurar un currículum de manera adecuada, evitar errores de lectura al procesar el archivo y asegurar que las habilidades y proyectos del candidato lleguen de forma clara a los reclutadores y evaluadores técnicos.

\subsection{Portales de empresas Big Tech y conglomerados tecnológicos}
\label{sec:portalesCorporativos}

Las grandes compañías tecnológicas y las empresas del grupo FAANG (Meta, Apple, Amazon, Netflix y Google, junto con otros líderes de la industria) concentran buena parte de las convocatorias más buscadas de pasantías en ingeniería de software. Debido al gran volumen de solicitudes que reciben a nivel internacional, estas empresas utilizan portales propietarios o plataformas corporativas adaptadas a sus necesidades, desde donde gestionan las etapas iniciales de filtrado (\textit{screening}) y el envío de evaluaciones de programación en línea (\textit{online assessment}, OA).

A partir de las convocatorias registradas en el repositorio \textit{Summer 2027 Tech Internships} \parencite{simplifyjobs2026summer}, se presenta la lista de vacantes activas correspondientes a este grupo corporativo, detallando el puesto, la plataforma utilizada y el enlace directo para enviar la solicitud:

\begin{itemize}[leftmargin=0.2in, itemsep=0.3em]
    \item \textbf{Google:} Portal corporativo propio \parencite{google2026careers}.
    \begin{itemize}[leftmargin=0.15in, noitemsep]
        \item \textit{Puesto:} Forward Deployed Engineer Intern.
        \item \textit{Enlace directo de aplicación:}\\ {\small\url{https://www.google.com/about/careers/applications/jobs/results/75665705386549958}}
    \end{itemize}

    \item \textbf{Amazon:} Portal corporativo propio \parencite{amazon2026jobs}.
    \begin{itemize}[leftmargin=0.15in, noitemsep]
        \item \textit{Puesto:} Software Development Engineer Intern (Machine Learning / Automated Reasoning).
        \item \textit{Enlace directo de aplicación:}\\ {\small\url{https://amazon.jobs/en/jobs/10593521/software-development-engineer-intern-machine-learning-automated-reasoning-summer-2027}}
    \end{itemize}

    \item \textbf{NVIDIA:} Portal corporativo en Workday HCM.
    \begin{itemize}[leftmargin=0.15in, noitemsep]
        \item \textit{Puesto:} AI Systems Software Engineer Intern.
        \item \textit{Enlace directo de aplicación:}\\ {\small\url{https://nvidia.wd5.myworkdayjobs.com/NVIDIAExternalCareerSite/job/US-TX-Austin/AI-Systems-Software-Engineering-Intern---Summer-2027_JR2027170}}
    \end{itemize}

    \item \textbf{Netflix:} Portal corporativo propio.
    \begin{itemize}[leftmargin=0.15in, noitemsep]
        \item \textit{Puesto:} Software Engineer Intern.
        \item \textit{Enlace directo de aplicación:}\\ {\small\url{https://explore.jobs.netflix.net/careers/job/790318904468}}
    \end{itemize}

    \item \textbf{Tesla:} Portal corporativo propio.
    \begin{itemize}[leftmargin=0.15in, noitemsep]
        \item \textit{Puesto:} Internship - Software Engineering (People Products).
        \item \textit{Enlace directo de aplicación:}\\ {\small\url{https://www.tesla.com/careers/search/job/284004}}
    \end{itemize}

    \item \textbf{LinkedIn:} Portal de empleo en SmartRecruiters.
    \begin{itemize}[leftmargin=0.15in, noitemsep]
        \item \textit{Puesto:} Software Engineer Intern.
        \item \textit{Enlace directo de aplicación:}\\ {\small\url{https://jobs.smartrecruiters.com/LinkedIn3/744000156616169}}
    \end{itemize}

    \item \textbf{Salesforce:} Portal Futureforce en Workday HCM.
    \begin{itemize}[leftmargin=0.15in, noitemsep]
        \item \textit{Puesto:} Summer 2027 Intern - Software Engineer (Security \& Government).
        \item \textit{Enlace directo de aplicación:}\\ {\footnotesize\url{https://salesforce.wd12.myworkdayjobs.com/External_Career_Site/job/California---San-Francisco/Summer-2027-Intern---Software-Engineer--Security-Government-_JR364911-1}}
    \end{itemize}

    \item \textbf{Stripe:} Tablero de empleo en Greenhouse Software.
    \begin{itemize}[leftmargin=0.15in, noitemsep]
        \item \textit{Puesto:} Software Engineer Intern.
        \item \textit{Enlace directo de aplicación:}\\ {\small\url{https://stripe.com/jobs/search?gh_jid=8241260}}
    \end{itemize}

    \item \textbf{Figma:} Tablero de empleo en Greenhouse Software.
    \begin{itemize}[leftmargin=0.15in, noitemsep]
        \item \textit{Puesto:} Software Engineer Intern.
        \item \textit{Enlace directo de aplicación:}\\ {\small\url{https://boards.greenhouse.io/figma/jobs/6143238004}}
    \end{itemize}

    \item \textbf{Datadog:} Tablero de empleo en Greenhouse Software.
    \begin{itemize}[leftmargin=0.15in, noitemsep]
        \item \textit{Puesto:} Software Engineer Intern - Summer.
        \item \textit{Enlace directo de aplicación:}\\ {\small\url{https://careers.datadoghq.com/detail/8052118/?gh_jid=8052118}}
    \end{itemize}

    \item \textbf{Databricks:} Tablero de empleo en Greenhouse Software.
    \begin{itemize}[leftmargin=0.15in, noitemsep]
        \item \textit{Puesto:} Software Engineer Intern.
        \item \textit{Enlace directo de aplicación:}\\ {\small\url{https://boards.greenhouse.io/embed/job_app?token=6865687002}}
    \end{itemize}

    \item \textbf{Cloudflare:} Tablero de empleo en Greenhouse Software.
    \begin{itemize}[leftmargin=0.15in, noitemsep]
        \item \textit{Puesto:} Software Engineer Intern.
        \item \textit{Enlace directo de aplicación:}\\ {\small\url{https://boards.greenhouse.io/cloudflare/jobs/8199958}}
    \end{itemize}

    \item \textbf{Coinbase:} Tablero de empleo en Greenhouse Software.
    \begin{itemize}[leftmargin=0.15in, noitemsep]
        \item \textit{Puesto:} Software Engineer Intern.
        \item \textit{Enlace directo de aplicación:}\\ {\small\url{https://boards.greenhouse.io/embed/job_app?token=8168315}}
    \end{itemize}

    \item \textbf{Waymo:} Tablero de empleo en Greenhouse Software.
    \begin{itemize}[leftmargin=0.15in, noitemsep]
        \item \textit{Puesto:} AI-driven ML Performance Engineering Intern.
        \item \textit{Enlace directo de aplicación:}\\ {\small\url{https://careers.withwaymo.com/jobs?gh_jid=8248060}}
    \end{itemize}

    \item \textbf{Lyft:} Tablero de empleo en Greenhouse Software (Careerpuck).
    \begin{itemize}[leftmargin=0.15in, noitemsep]
        \item \textit{Puesto:} Software Engineer Intern (Backend).
        \item \textit{Enlace directo de aplicación:}\\ {\small\url{https://app.careerpuck.com/job-board/lyft/job/8767726002?gh_jid=8767726002}}
    \end{itemize}

    \item \textbf{Adobe:} Portal corporativo en Workday HCM.
    \begin{itemize}[leftmargin=0.15in, noitemsep]
        \item \textit{Puesto:} Intern - AI \& Architecture.
        \item \textit{Enlace directo de aplicación:}\\ {\footnotesize\url{https://adobe.wd5.myworkdayjobs.com/external_experienced/job/San-Jose/XMLNAME-2027-Intern---AI---Architecture_R172998}}
    \end{itemize}
\end{itemize}

\subsection{Sistemas de seguimiento de candidatos}
\label{sec:sistemasSeguimiento}

Además de los portales propios de las grandes corporaciones, la mayoría de las empresas tecnológicas medianas, empresas emergentes (\textit{startups}) y firmas especializadas en ciberseguridad o infraestructura utilizan plataformas comerciales en la nube bajo el modelo de software como servicio (\textit{software as a service}, SaaS). Estas herramientas estandarizan la recepción de candidaturas, automatizan el filtrado inicial y conectan a los aspirantes con los equipos de contratación:

\begin{itemize}[leftmargin=0.2in, itemsep=0.4em]
    \item \textbf{Greenhouse Software:} Es uno de los sistemas ATS más extendidos entre empresas de tecnología y ciberseguridad (como KnowBe4, Discord o Stripe). Sus tableros públicos de empleo se alojan bajo el dominio \url{https://job-boards.greenhouse.io/} \parencite{greenhouse2026software}, con formularios directos donde se carga el currículum en formato PDF y se ingresan enlaces a perfiles técnicos en GitHub y LinkedIn.

    \item \textbf{Workday Human Capital Management:} Es la plataforma habitual en grandes corporaciones de sectores industriales, tecnológicos y de salud (como PTC, VSP Vision, NVIDIA y Adobe). Su módulo de reclutamiento \parencite{workday2026hcm} funciona con subdominios específicos por empresa (\url{https://empresa.wd1.myworkdayjobs.com/}), donde el aspirante crea una cuenta y captura su historial académico y laboral paso a paso.

    \item \textbf{Lever:} Utilizada por diversas empresas en crecimiento, Lever \parencite{lever2026recruiting} combina funciones de ATS con el seguimiento de candidatos en un formulario de una sola página, lo que agiliza el envío de solicitudes al reducir los pasos intermedios.

    \item \textbf{Agregadores comunitarios y repositorios abiertos:} Para dar seguimiento a cientos de vacantes dispersas en múltiples portales, las comunidades de estudiantes de ingeniería en computación organizan listas abiertas en tiempo real. El referente más utilizado para pasantías estivales es el repositorio colaborativo del Pitt Computer Science Club y la organización Simplify: \textit{Summer 2027 Tech Internships} \parencite{simplifyjobs2026summer}, donde se concentran convocatorias verificadas con enlaces directos de postulación.
\end{itemize}

En el Cuadro~\ref{tab:comparativaPlataformas} se sintetiza el funcionamiento de las principales plataformas analizadas, su enfoque para procesar solicitudes y ejemplos de organizaciones que las utilizan.

\begin{table}[htbp]
    \centering
    \small
    \caption{Comparativa de plataformas y portales de reclutamiento en ingeniería de software.}
    \label{tab:comparativaPlataformas}
    \begin{tabularx}{\textwidth}{p{3.2cm} p{2.6cm} p{3.4cm} X}
        \toprule
        \textbf{Plataforma / Portal} & \textbf{Tipo de sistema} & \textbf{Procesamiento de datos} & \textbf{Empresas y casos de uso} \\
        \midrule
        Portales Big Tech y FAANG & Portales propios en la web & Lectura automática de currículum y asignación de pruebas técnicas (OA) & Google, Amazon, Netflix, Tesla, NVIDIA. Gestión de miles de aspirantes en todo el mundo. \\
        \addlinespace
        Amazon Jobs & Portal corporativo propio & Búsqueda por número de vacante y evaluación técnica estandarizada & Amazon y AWS. Convocatorias globales para desarrollo de software (SDE). \\
        \addlinespace
        Greenhouse Software & ATS en la nube (SaaS) & Formulario optimizado, lectura de PDF y campos para enlaces técnicos & KnowBe4, Discord, Stripe. Muy utilizado en empresas de software y ciberseguridad. \\
        \addlinespace
        Workday HCM & Plataforma empresarial (ERP / ATS) & Registro modular por etapas; requiere cuenta por empresa & PTC, VSP Vision, Adobe, NVIDIA. Grandes corporaciones multinacionales. \\
        \addlinespace
        Lever & ATS en la nube (SaaS) & Formulario directo en una sola página con carga de archivos y perfiles & Empresas medianas y \textit{startups} en crecimiento. \\
        \addlinespace
        SimplifyJobs / Repositorios abiertos & Agregador colaborativo en GitHub & Rastreo comunitario de vacantes con enlaces oficiales directos & Seguimiento en tiempo real de convocatorias de verano para estudiantes. \\
        \bottomrule
    \end{tabularx}
\end{table}

\section{Casos prácticos de envío de solicitudes}
\label{sec:casosPostulacion}

Llenar una solicitud para un puesto técnico de ingeniería requiere adaptar los datos académicos, los proyectos personales y la experiencia laboral al formato específico que solicita cada empresa. Dependiendo de la compañía y de la plataforma que utilice, el proceso puede requerir desde un formulario rápido de una página hasta el llenado detallado de módulos de educación, enlaces a repositorios públicos o respuestas a preguntas de filtro sobre disponibilidad y dominio de tecnologías.

A continuación se presentan cinco ejemplos prácticos de postulación a vacantes de pasantía de verano (\textit{internships}) tomadas del repositorio \textit{Summer 2027 Tech Internships} \parencite{simplifyjobs2026summer}. Para cada caso se describe el puesto, la empresa, la plataforma utilizada, el enlace de postulación y la organización del perfil del estudiante (formación en la UNAM, proyectos en GitHub, concursos de programación y trayectoria laboral previa).

\subsection{Caso 1: Solicitud en ByteDance y TikTok (Software Engineer Intern)}
\label{sec:casoByteDance}

ByteDance y su filial TikTok reciben las postulaciones técnicas mediante un portal de empleo corporativo unificado \parencite{bytedance2026jobs, tiktok2026careers}. La vacante seleccionada corresponde al puesto de pasantía en ingeniería de software:

\begin{itemize}[leftmargin=0.2in, itemsep=0.2em]
    \item \textbf{Empresa:} ByteDance Ltd. / TikTok Inc.
    \item \textbf{Puesto:} Software Engineer Intern.
    \item \textbf{Plataforma de postulación:} Portal corporativo global de empleo.
    \item \textbf{Enlaces directos de aplicación:}
    \begin{itemize}[leftmargin=0.15in, noitemsep]
        \item ByteDance: {\footnotesize\url{https://jobs.bytedance.com/en/position/7686394581777631541/detail}}
        \item TikTok: {\footnotesize\url{https://lifeattiktok.com/search/7701069870360627509}}
    \end{itemize}
\end{itemize}

La información del perfil se distribuye en las secciones del formulario de la siguiente manera:

\begin{enumerate}[leftmargin=0.25in, itemsep=0.35em]
    \item \textbf{Formación académica (\textit{Education}):}
    \begin{itemize}[leftmargin=0.15in, noitemsep]
        \item \textit{Institución:} Universidad Nacional Autónoma de México (UNAM), Facultad de Ingeniería.
        \item \textit{Carrera:} Licenciatura en Ingeniería en Computación (\textit{Bachelor's degree in Computer Engineering}).
        \item \textit{Periodo escolar:} Septiembre de 2024 a mayo de 2028 (fecha estimada de graduación).
        \item \textit{Materias relevantes:} Estructuras de Datos y Algoritmos, Programación Orientada a Objetos.
    \end{itemize}

    \item \textbf{Experiencia laboral y servicio social (\textit{Work \& Internship Experience}):}
    \begin{itemize}[leftmargin=0.15in, noitemsep]
        \item \textit{Ambar Electricians (Senior Installer, octubre de 2018 a la fecha):} Técnico líder en instalaciones eléctricas residenciales y comerciales y sistemas solares fotovoltaicos. Trabajo de campo enfocado en interpretación de planos eléctricos, cumplimiento de normas de seguridad, diagnóstico de fallas y entrega técnica de instalaciones funcionales.
        \item \textit{iOS Development Lab FI UNAM (Social Service Intern, febrero de 2026 a la fecha):} Programa de servicio social en la Facultad de Ingeniería enfocado en el desarrollo de aplicaciones móviles con impacto social en Swift y SwiftUI, aplicando la arquitectura modelo-vista-controlador (MVC) a partir de necesidades comunitarias detectadas en campo.
    \end{itemize}

    \item \textbf{Proyectos de desarrollo (\textit{Project Experience}):}
    \begin{itemize}[leftmargin=0.15in, noitemsep]
        \item \textit{Investment in Solar Systems (Sunstake):} Aplicación móvil para facilitar la inversión en energía solar eliminando la complejidad de la cadena de bloques. Utiliza un modelo de inteligencia artificial en el dispositivo para explicar contratos inteligentes en lenguaje sencillo y una billetera no custodial para resguardar las claves criptográficas (\url{https://github.com/Oscar-Cabr/Sunstake-HackApp}).
        \item \textit{School Management System:} Sistema de control escolar desarrollado en Java bajo principios de programación orientada a objetos (abstracción, herencia, polimorfismo y encapsulamiento), implementando los patrones de diseño \textit{Singleton} y \textit{Memento} para ordenar la creación de objetos y el control de estado (\url{https://github.com/Oscar-Cabr/Proyecto_Gestion_Escolar}).
    \end{itemize}

    \item \textbf{Muestras de código y perfiles profesionales (\textit{Work Samples \& SNS}):}
    \begin{itemize}[leftmargin=0.15in, noitemsep]
        \item \textit{Repositorio de Sunstake:}\\ {\footnotesize\url{https://github.com/Oscar-Cabr/Sunstake-HackApp}}
        \item \textit{Perfil de programación competitiva (Codeforces):}\\ {\footnotesize\url{https://codeforces.com/profile/RacsoFractal}} (358 problemas resueltos en C++).
        \item \textit{Sitio de clases de CPCFI:}\\ {\footnotesize\url{https://oscar-cabr.github.io/CPCFI-Classes/}}
        \item \textit{Perfiles profesionales:}
        \begin{itemize}[leftmargin=0.15in, noitemsep]
            \item LinkedIn: {\footnotesize\url{https://www.linkedin.com/in/oscar-cabrera-1698b0289/}}
            \item GitHub: {\footnotesize\url{https://github.com/Oscar-Cabr}}
        \end{itemize}
    \end{itemize}

    \item \textbf{Reconocimientos y concursos (\textit{Honors \& Awards}):}
    \begin{itemize}[leftmargin=0.15in, noitemsep]
        \item \textit{ICPC Subregional (2025):} Clasificación en el \textit{top} 70 de más de 500 equipos en la región México, resolviendo problemas algorítmicos en C++.
        \item \textit{IEEEXtreme 19.0 (2025):} Clasificación en el \textit{top} 20 de 130 equipos en la sección México durante la competencia internacional de 24 horas.
        \item \textit{Finalista nacional en Swift Changemakers Hackathon (2026):} Clasificación en el \textit{top} 3 local y pase a la fase nacional del certamen de Enactus México en la categoría de energía.
    \end{itemize}

    \item \textbf{Habilidades técnicas e idiomas (\textit{Skills \& Languages}):}
    \begin{itemize}[leftmargin=0.15in, noitemsep]
        \item \textit{Lenguajes y herramientas:} C, C++, Swift, Java, Python; Linux (Arch, Debian), Git, Docker, Neovim, Shell Scripting, LaTeX.
        \item \textit{Idiomas:} Español (nativo), Inglés (profesional avanzado, C1), Chino Mandarín (elemental, HSK1).
    \end{itemize}

    \item \textbf{Presentación personal en el portal (\textit{Self-Introduction}):}
    Texto ingresado en el campo libre del formulario:
    \begin{quotation}\small\upshape
    \enquote{I'm a Computer Engineering student at UNAM's Facultad de Ingeniería, with a strong foundation in low-level programming and algorithms. C and C++ are my core languages, and I use them to solve problems under competitive constraints: I compete in the ICPC subregional, where my team placed in the top 70 out of 500+ teams in the Mexico region, and in IEEEXtreme 18.0 and 19.0, placing in the top 20 out of 130 teams in the Mexico section.

    I'm also the general coordinator of CPCFI, my faculty's competitive programming club, where I lead the academic team. I've mentored 50+ competitors and written the lecture material for more than ten advanced topics, delivered across 30+ technical sessions. I direct the design of our contests and author original problems for them, and I coordinate events with other schools and competitive programming organizations --- most recently running both phases and both divisions of the 2025--2026 Interpreparatoriano Programming Contest for UNAM's Escuela Nacional Preparatoria, including the on-site venue for the final.

    On the development side, I'm a social service intern at the iOS Development Lab, building social-impact iOS applications in Swift and SwiftUI. That work led to Sunstake, a solar energy investment app I built for the Swift Changemakers Hackathon organized by Enactus México, where my team advanced to the national finals. The domain wasn't arbitrary: I've worked since 2018 as a senior installer at Ambar Electricians, handling residential and commercial electrical and solar installations --- experience that taught me to design around real-world constraints and around the person who ultimately has to use the system.

    I'm looking for a software engineering internship where I can take on demanding technical problems and learn from experienced engineers.}
    \end{quotation}
\end{enumerate}

\subsection{Caso 2: Solicitud en Amazon (Software Development Engineer Intern)}
\label{sec:casoAmazon}

Amazon centraliza sus convocatorias de verano a través de su portal propio Amazon Jobs \parencite{amazon2026jobs}. La postulación para la vacante de desarrollo de software enfocada en aprendizaje automático y razonamiento automatizado se organiza de la siguiente manera:

\begin{itemize}[leftmargin=0.2in, itemsep=0.2em]
    \item \textbf{Empresa:} Amazon.com, Inc. / AWS.
    \item \textbf{Puesto:} Software Development Engineer Intern (Machine Learning / Automated Reasoning).
    \item \textbf{Plataforma de postulación:} Amazon Jobs Portal.
    \item \textbf{Enlace directo de aplicación:}\\ {\scriptsize\url{https://amazon.jobs/en/jobs/10593521/software-development-engineer-intern-machine-learning-automated-reasoning-summer-2027}}
\end{itemize}

El perfil del candidato se orienta para responder a las expectativas de la vacante:
\begin{itemize}[leftmargin=0.2in, itemsep=0.3em]
    \item \textbf{Enfoque algorítmico:} Se resalta la participación en el ICPC e IEEEXtreme y la práctica continua en C++ como prueba de habilidad para analizar problemas complejos y optimizar algoritmos bajo límites estrictos de tiempo y memoria.
    \item \textbf{Proyectos técnicos:} Se destaca el proyecto Sunstake para mostrar experiencia práctica integrando modelos de inteligencia artificial en aplicaciones móviles, junto con el material educativo de CPCFI sobre estructuras de datos y análisis de complejidad.
    \item \textbf{Principios de liderazgo de Amazon (\textit{Leadership Principles}):} Se vincula la experiencia de diagnóstico de fallas en Ambar Electricians y la coordinación de cursos en el club de programación con principios de Amazon como profundizar en los temas (\textit{Dive Deep}), tomar la iniciativa (\textit{Bias for Action}) y dar resultados (\textit{Deliver Results}).
\end{itemize}

\subsection{Caso 3: Solicitud en KnowBe4 (Software Engineer Intern)}
\label{sec:casoKnowBe4}

KnowBe4, empresa especializada en ciberseguridad y capacitación contra ataques de ingeniería social, recibe sus solicitudes a través de Greenhouse Software \parencite{greenhouse2026software}:

\begin{itemize}[leftmargin=0.2in, itemsep=0.2em]
    \item \textbf{Empresa:} KnowBe4, Inc.
    \item \textbf{Puesto:} Software Engineer Intern.
    \item \textbf{Plataforma de postulación:} Tablero de empleo en Greenhouse Software.
    \item \textbf{Enlace directo de aplicación:}\\ {\footnotesize\url{https://job-boards.greenhouse.io/knowbe4/jobs/8919280002}}
\end{itemize}

El llenado de la postulación en Greenhouse se estructura considerando los siguientes puntos:
\begin{itemize}[leftmargin=0.2in, itemsep=0.3em]
    \item \textbf{Formato del currículum para el ATS:} Se adjunta el archivo en formato PDF limpio, con texto claro y sin columnas complicadas ni tablas anidadas, asegurando que el lector automático de Greenhouse reconozca sin problemas las tecnologías del perfil (C, C++, Java, Linux, Docker y Git).
    \item \textbf{Seguridad y buenas prácticas de desarrollo:} Siendo una empresa de seguridad, se destacan las medidas de protección implementadas en Sunstake (manejo de claves criptográficas en billeteras no custodiales) y el diseño estructurado del sistema escolar en Java con encapsulamiento estricto.
    \item \textbf{Preguntas de filtro del formulario:} Se completan los campos sobre disponibilidad a tiempo completo para el verano de 2027 y la condición de estudiante activo en la carrera de Ingeniería en Computación de la UNAM.
\end{itemize}

\subsection{Caso 4: Solicitud en VSP Vision (AI \& Machine Learning Intern)}
\label{sec:casoVspVision}

VSP Vision gestiona sus contrataciones de pasantías técnicas en salud visual y tecnologías de inteligencia artificial a través de la plataforma empresarial Workday HCM \parencite{workday2026hcm}:

\begin{itemize}[leftmargin=0.2in, itemsep=0.2em]
    \item \textbf{Empresa:} VSP Vision.
    \item \textbf{Puesto:} AI, Agentic \& Machine Learning Development Intern.
    \item \textbf{Plataforma de postulación:} Portal corporativo en Workday HCM.
    \item \textbf{Enlace directo de aplicación:}\\ {\scriptsize\url{https://vsp.wd1.myworkdayjobs.com/vspvisioncareers/job/Remote-US/Internship--AI---Agentic---Machine-Learning-Development_R-10072}}
\end{itemize}

Los datos del candidato se capturan en las diferentes etapas del formulario de Workday:
\begin{itemize}[leftmargin=0.2in, itemsep=0.3em]
    \item \textbf{Registro por módulos:} Se capturan paso a paso los estudios de ingeniería en computación en la UNAM, el servicio social en el iOS Development Lab y la experiencia técnica como instalador en Ambar Electricians.
    \item \textbf{Proyectos de inteligencia artificial aplicada:} Se presenta la aplicación Sunstake destacando el uso de modelos de lenguaje en el dispositivo para explicar operaciones financieras de manera comprensible al usuario, siguiendo principios de diseño centrado en las personas.
    \item \textbf{Bases lógicas y de arquitectura:} Se demuestra habilidad para estructurar flujos de software claros mediante patrones de diseño (como el patrón \textit{Memento} en el gestor escolar) y algoritmos bien fundamentados desarrollados en las sesiones de CPCFI.
\end{itemize}

\subsection{Caso 5: Solicitud en PTC (Software Engineering Intern)}
\label{sec:casoPtc}

PTC Inc., empresa dedicada al desarrollo de software para diseño asistido por computadora (CAD), gestión del ciclo de vida de productos (PLM) e Internet de las cosas (\textit{Internet of Things}, IoT), procesa sus vacantes mediante Workday HCM \parencite{workday2026hcm}:

\begin{itemize}[leftmargin=0.2in, itemsep=0.2em]
    \item \textbf{Empresa:} PTC Inc.
    \item \textbf{Puesto:} Software Engineering Intern.
    \item \textbf{Plataforma de postulación:} Portal corporativo en Workday HCM.
    \item \textbf{Enlace directo de aplicación:}\\ {\scriptsize\url{https://ptc.wd1.myworkdayjobs.com/ptc/job/Cambridge-United-Kingdom/Software-Engineering-Summer-Internship-2027---Cambridge_JR112998}}
\end{itemize}

El perfil del estudiante se conecta con las necesidades de desarrollo de PTC:
\begin{itemize}[leftmargin=0.2in, itemsep=0.3em]
    \item \textbf{Combinación de experiencia física y desarrollo de software:} PTC crea productos que interactúan con equipos físicos e industriales. Aquí resulta muy útil combinar la experiencia de campo en Ambar Electricians (interpretación de diagramas eléctricos, normas de seguridad y puesta en marcha de sistemas fotovoltaicos) con los conocimientos de programación.
    \item \textbf{Programación de bajo nivel y código orientado a objetos:} Se aportan muestras de código en C y C++ enfocadas en rendimiento y manejo eficiente de memoria, junto con el proyecto de control escolar en Java como muestra de código modular y bien organizado.
    \item \textbf{Liderazgo y trabajo en equipo:} Se respalda la capacidad para colaborar en equipos multidisciplinarios a través de la coordinación académica del club CPCFI y la tutoría impartida a más de 50 compañeros.
\end{itemize}

\section{Conclusiones personales}
\label{sec:conclusionesPersonales}

Analizar cómo funcionan las plataformas de contratación en las empresas de tecnología y preparar solicitudes reales deja una lección muy clara: buscar trabajo o pasantías en la industria no es algo que deba dejarse para el último semestre de la carrera. Las convocatorias para pasantías de verano abren prácticamente con un año de anticipación, como se ve con las ofertas para 2027 disponibles desde 2026. Quien espera a terminar las materias para empezar a armar su currículum o practicar para las entrevistas técnicas suele encontrarse con que los filtros iniciales, tanto los sistemas automáticos de lectura como las pruebas de código en línea, exigen un nivel de preparación que no se improvisa en un par de semanas. Prepararse con tiempo, entendiendo qué buscan los reclutadores y cómo evalúan los equipos de ingeniería, permite llegar a estos procesos con seguridad y con un perfil competitivo.

En este sentido, la vida universitaria va mucho más allá de lo que pasa dentro del salón de clases. Las materias de la carrera dan las bases teóricas indispensables en matemáticas, lógica, sistemas operativos y estructuras de datos; sin embargo, es en los proyectos y actividades extracurriculares donde ese conocimiento realmente se pone a prueba. Participar activamente en clubes estudiantiles, como el Club de Programación Competitiva de la Facultad de Ingeniería (CPCFI), y competir en torneos como el ICPC o el certamen IEEEXtreme acostumbra al estudiante a resolver problemas difíciles bajo presión de tiempo, a escribir código limpio y eficiente en lenguajes como C y C++, y a pensar en la complejidad de sus soluciones. Además, coordinar grupos de estudio, dar clases a otros compañeros y diseñar problemas técnicos ayuda a desarrollar la habilidad de comunicar ideas con claridad, algo fundamental para defender una propuesta técnica en una entrevista de diseño o explicar el funcionamiento de un sistema.

Por otra parte, tener proyectos propios y código público en plataformas como GitHub marca una gran diferencia al momento de postularse. Decir en un currículum que se sabe programar en un lenguaje o que se conocen ciertos patrones de diseño no dice mucho si no hay evidencia visible de ello. Desarrollar aplicaciones reales, ya sea a través de iniciativas comunitarias como el servicio social en el iOS Development Lab o creando proyectos independientes como la aplicación Sunstake, obliga a enfrentarse a problemas concretos: diseñar arquitecturas mantenibles, manejar el estado de una aplicación, integrar herramientas modernas como modelos de inteligencia artificial en el dispositivo y cuidar la experiencia del usuario. Todo esto demuestra a los reclutadores que uno no se queda solo con lo que pide el profesor en una práctica escolar, sino que tiene la iniciativa de investigar, construir soluciones completas y aplicar buenas prácticas de ingeniería de software.

Finalmente, cualquier experiencia laboral o práctica que se haya tenido durante los años de estudio, aun cuando no pertenezca estrictamente al desarrollo de software de oficina, aporta un valor enorme si se sabe enfocar. Por ejemplo, trabajar durante años en instalaciones eléctricas y proyectos solares enseña a interpretar planos técnicos, a diagnosticar fallas en sistemas reales bajo normas estrictas de seguridad y a asumir la responsabilidad técnica de un trabajo de principio a fin. Esa visión práctica de campo complementa el pensamiento lógico de la computación, recordando que los programas no existen en el vacío, sino que corren en computadoras reales y resuelven problemas para personas reales. Al final, prepararse desde la universidad para la vida laboral consiste en combinar esa base teórica con la práctica constante, aprovechando cada proyecto escolar, cada concurso y cada experiencia de trabajo para construir un perfil técnico sólido y honesto.

\nocite{*}
\printbibliography[heading=bibintoc, title={Referencias}]

\end{document}
